% Zone „Das kennst du schon“ – aus bank/brueche-dezimalzahlen/zone.jsonl (zusammenbau v0.1)
\einheitenkopf[zone][Kennst du schon]{Das kennst du schon}
\setcounter{aufgabe}{0}
%% TODO Grundvorstellungs-Aufgabe als erste Hauptnummer der Zone (2.8) fehlt in der Bank

%% TODO Titel: Ich-kann-Satz fehlt in der Bank (Platzhalter: Kettenname) – Nr. 1
%% TODO Anweisung: ein Satz über der ersten Teilaufgabe fehlt in der Bank – Nr. 1
\begin{aufgabe}{Teilen und Vervielfachen im Kopf (vierundzwanzig geteilt durch sechs, drei mal neun), Einmaleins}
%% TODO Darstellung für schwach fehlt in der Bank (brueche-dezimalzahlen-zone-f1-v1; Abschnitt „Für schwache Schüler“)
\swz{$56 : 7$}{}
%% TODO Darstellung für schwach fehlt in der Bank (brueche-dezimalzahlen-zone-f1-v3; Abschnitt „Für schwache Schüler“)
\swz{$4{,}2 : 6$}{}
\end{aufgabe}

%% TODO Titel: Ich-kann-Satz fehlt in der Bank (Platzhalter: Kettenname) – Nr. 2
%% TODO Anweisung: ein Satz über der ersten Teilaufgabe fehlt in der Bank – Nr. 2
\begin{aufgabe}{Größen mit Komma lesen und umrechnen (1,2 kg = 1200 g; 2,50 €)}
%% TODO Darstellung für schwach fehlt in der Bank (brueche-dezimalzahlen-zone-f2-v1; Abschnitt „Für schwache Schüler“)
\swz{$3{,}5$ kg – wie viel Gramm?}{}
%% TODO Darstellung für schwach fehlt in der Bank (brueche-dezimalzahlen-zone-f2-v3; Abschnitt „Für schwache Schüler“)
\swz{$0{,}85$ l – wie viel Milliliter?}{}
\end{aufgabe}

%% TODO Titel: Ich-kann-Satz fehlt in der Bank (Platzhalter: Kettenname) – Nr. 3
%% TODO Anweisung: ein Satz über der ersten Teilaufgabe fehlt in der Bank – Nr. 3
\begin{aufgabe}{Kreissektor als Anteil vom Vollkreis (120° von 360°)}
%% TODO Darstellung für schwach fehlt in der Bank (brueche-dezimalzahlen-zone-f3-v1; Abschnitt „Für schwache Schüler“)
\swz{$180^\circ$ von $360^\circ$ – welcher Anteil vom Vollkreis?}{}
%% TODO Darstellung für schwach fehlt in der Bank (brueche-dezimalzahlen-zone-f3-v3; Abschnitt „Für schwache Schüler“)
\swz{$45^\circ$ von $360^\circ$ – welcher Anteil vom Vollkreis?}{}
\end{aufgabe}

%% TODO Titel: Ich-kann-Satz fehlt in der Bank (Platzhalter: Kettenname) – Nr. 4
%% TODO Anweisung: ein Satz über der ersten Teilaufgabe fehlt in der Bank – Nr. 4
\begin{aufgabe}{Teiler und Vielfache einer Zahl angeben (Teiler von zwölf; Vielfache von vier)}
%% TODO Darstellung für schwach fehlt in der Bank (brueche-dezimalzahlen-zone-f4-v1; Abschnitt „Für schwache Schüler“)
\swz{Nenne alle Teiler von $10$.}{}
%% TODO Darstellung für schwach fehlt in der Bank (brueche-dezimalzahlen-zone-f4-v3; Abschnitt „Für schwache Schüler“)
\swz{Nenne alle Teiler von $36$.}{}
\end{aufgabe}

%% TODO Titel: Ich-kann-Satz fehlt in der Bank (Platzhalter: Kettenname) – Nr. 5
%% TODO Anweisung: ein Satz über der ersten Teilaufgabe fehlt in der Bank – Nr. 5
\begin{aufgabe}{Stellenwerte natürlicher Zahlen und Zahlenstrahl mit natürlichen Zahlen (Skala ablesen)}
%% TODO Darstellung für schwach fehlt in der Bank (brueche-dezimalzahlen-zone-f5-v1; Abschnitt „Für schwache Schüler“)
\swz{Welchen Stellenwert hat die $7$ in $4\,782$?}{}
%% TODO Darstellung für schwach fehlt in der Bank (brueche-dezimalzahlen-zone-f5-v3; Abschnitt „Für schwache Schüler“)
\swz{Von $0$ bis $1\,000$ sind $10$ gleich lange Abschnitte. Welche Zahl gehört an den dritten Teilstrich nach der $0$?}{}
\end{aufgabe}

%% TODO Titel: Ich-kann-Satz fehlt in der Bank (Platzhalter: Kettenname) – Nr. 6
%% TODO Anweisung: ein Satz über der ersten Teilaufgabe fehlt in der Bank – Nr. 6
\begin{aufgabe}{Schriftlich oder mit Taschenrechner dividieren (drei geteilt durch acht)}
%% TODO Darstellung für schwach fehlt in der Bank (brueche-dezimalzahlen-zone-f6-v1; Abschnitt „Für schwache Schüler“)
\swz{$84 : 4$}{}
%% TODO Darstellung für schwach fehlt in der Bank (brueche-dezimalzahlen-zone-f6-v3; Abschnitt „Für schwache Schüler“)
\swz{$7 : 4$}{}
\end{aufgabe}

%% TODO Titel: Ich-kann-Satz fehlt in der Bank (Platzhalter: Kettenname) – Nr. 7
%% TODO Anweisung: ein Satz über der ersten Teilaufgabe fehlt in der Bank – Nr. 7
\begin{aufgabe}{Prozent als Hundertstel (45 \% = 0,45)}
%% TODO Darstellung für schwach fehlt in der Bank (brueche-dezimalzahlen-zone-f7-v1; Abschnitt „Für schwache Schüler“)
\swz{$30\,\%$ – als Bruch mit dem Nenner $100$?}{}
%% TODO Darstellung für schwach fehlt in der Bank (brueche-dezimalzahlen-zone-f7-v3; Abschnitt „Für schwache Schüler“)
\swz{$1{,}2$ – wie viel Prozent?}{}
\end{aufgabe}

%% TODO Titel: Ich-kann-Satz fehlt in der Bank (Platzhalter: Kettenname) – Nr. 8
%% TODO Anweisung: ein Satz über der ersten Teilaufgabe fehlt in der Bank – Nr. 8
\begin{aufgabe}{Prozent als Hundertstel (45 \% = 0,45) – Fehler finden}
\swz[3]{Lea schreibt: $6\,\% = 0{,}6$. Finde den Fehler.}{}
\end{aufgabe}

%% TODO Titel: Ich-kann-Satz fehlt in der Bank (Platzhalter: Kettenname) – Nr. 9
%% TODO Anweisung: ein Satz über der ersten Teilaufgabe fehlt in der Bank – Nr. 9
\begin{aufgabe}{Prozent als Hundertstel (45 \% = 0,45) – selbst rechnen}
%% TODO Darstellung für schwach fehlt in der Bank (brueche-dezimalzahlen-zone-f7-v6; Abschnitt „Für schwache Schüler“)
\swz{$9\,\%$ – als Kommazahl?}{}
\end{aufgabe}
